\documentclass[12pt]{article}
\usepackage[margin=1in]{geometry}
\usepackage{enumitem}
\usepackage{titlesec}
\usepackage{parskip}
\usepackage{graphicx}
\usepackage{xcolor}
\usepackage{hyperref}
\usepackage{fontenc}

\title{\textbf{Lijphart 1991 RG Notes}\\
\large Lijphart, Arend. ``Constitutional Choices for New Democracies'' [1991].\\
\textit{Journal of Democracy} 2(1): 72-84.}
\author{}
\date{}

\begin{document}

\maketitle

Strongly supports combination of parliamentarian and proportional representation, particularly for new(er) democracies

Democratic quality: ``the degree to which a system meets such democratic norms as representativeness, accountability, equality, and participation''

``Proportionalists tend to attach greater importance to the representativeness of government, while plurality advocates view the capacity to govern as the more vital consideration''

Cites Powell's metrics (positively!). Also cites Dahl's rankings/methods/indicators of democratic quality

Generally supports moderate PR over extreme --- that is, some barriers to entry and minimal thresholds for representation

Very simple, pretty contrived article --- not a ton of evidence or particularly impressive model-building, but to a large extent, that as already been done --- this ends up more just being a restatement of existing arguments, plus a mild differentiation of moderate vs extreme PR

\end{document}
